\section{Limitations}
\label{sec:limits}
\emph{Deployment.} All measurements come from one host. The managed state is a local SQLite database, the governed operation is an ML-KEM-768 encapsulation round trip, and profile names are fixture labels, not evidence of protocol negotiation. The parameter sets ML-DSA-65 and ML-KEM-768 do not meet CNSA~2.0, which calls for Level~V parameters~\cite{cnsa2}, that is, ML-DSA-87 and ML-KEM-1024 (security category~5)~\cite{fips203,fips204}; costs were not measured for the Level~V sets. Keys are not hardware-protected, no validated module is involved, and liboqs is research software~\cite{liboqs}. All roles run in one process, so role keys separate provenance, not authority; a compromised process could use any key. The observers read the executor's store read-only: their observations are not derived from receipts but are not administratively independent of the executor. The witness keeps its checkpoints in the memory of that process, and no independent evidence retention is implemented. Durability was tested only with injected faults (crashes before and after the apply commit, failure of the controlled operation, restart by reopening the store), not with termination by the operating system, power loss, concurrent executors, or storage failure.

\emph{Observation.} The contract's outcome observer measures the recorded configuration, not delivered protection. The TLS panel is a standalone loopback experiment, not the contract's outcome observer; no transaction was closed on wire-level evidence. The conjunct $\Compatible$ checks the declared capabilities of listed endpoints; by Han's argument~\cite{han} it cannot establish what peers of other classes receive, and the panel covers only its four classes.

\emph{Advisor evaluation.} One model was evaluated, through \AdvInterface\ rather than the bare model endpoint; sampling parameters could not be set. Scenarios are synthetic, attacks are single-shot and non-adaptive, and each scenario was presented in one placement only, so placement is confounded with the scenario. Repetitions of a scenario are not independent, so counts are reported per trial and per scenario without pooled confidence intervals, and zero counts do not establish zero rates. The ablation removes the prompt sentences and the trust labels together and does not separate their effects. Because the model proposed no policy violation, containment rests on the scripted compromised proposer.

\emph{Zero-knowledge instantiation.} By the vendor's soundness calculator~\cite{risc0}, the worst accepted receipt has \ZkEndBits~bits of soundness under the toy-model conjecture, \ZkStrictBits~bits under the stricter conjecture, and \ZkProvenBits~bits proven (Section~\ref{sec:zklimits}). Zero knowledge, which the vendor has not proven~\cite{risc0sec}, and commitment hiding are assumed; a receipt reveals the padded execution length, and results, timing, and rationale can leak policy information. The program identifier is reproducible only with the same guest toolchain; reproduction on another machine was not evaluated.

\emph{Assurance.} The tests exercise the implemented checks on constructed fixtures; they neither estimate detection rates nor show that the obligations are complete. The contract and obligations lack a formal model and mechanical proof, and no assessor study was performed. Verifier memory grows linearly with transactions; archives beyond $10^5$ transactions were not evaluated. Policy adequacy, inventory completeness, and role-key protection remain assumptions.

\section{Conclusion}
An AI recommendation neither authorizes a PQC migration nor establishes that it occurred, and explaining a decision based on a confidential rule discloses the rule. Auditable by Design therefore binds audit evidence to the governed operation and controls disclosure of decision logic separately. Its transaction contract links authenticated observations, an exact-action approval consumed once, CAS execution, a read-only outcome observation, and a witnessed, ML-DSA-signed history; its verifier rejects correctly signed evidence that violates the cross-statement obligations.

Under complete proposer compromise, \AdvCompViolExecuted\ policy-violating proposals and \AdvCompInPolicyExecuted\ within-policy manipulations executed. With an LLM advisor and an approver accepting every gate-permitted proposal, within-policy manipulations executed in \AdvHardInPolicyExecuted\ trials under a hardened prompt and \AdvNaiveInPolicyExecuted\ under a naive one. Enforcement thus contained the tested policy violations independently of the proposer, whereas permitted choices depend on the advisor, prompt, and approver. Median costs on one host were \TimDurTotal~ms per governed transaction at steady state (mostly durable commits), \ScalHundredKMedR~s to verify a $10^5$-transaction archive, and \ZkProveMed~s to prove the confidential-predicate relation under conjectured soundness and assumed zero knowledge. A four-client panel detected two silent TLS downgrades, each visible to only one client class. Administrative role separation, wire-level observation inside the contract, an assessor study, and a formal model remain open.

\section*{Acknowledgment}
The author conceived this work, designed the approach, wrote the original manuscript draft, the initial implementation, and the initial benchmarks, and directed all subsequent work. Claude (Anthropic), used through Claude Code under the author's direction, revised, debugged, and extended the implementation and the benchmarks, added the LLM advisor harness, the RISC Zero guest and host, and the TLS observation module, ran the experiments, and restructured and edited the author's draft into the present version. The author reviewed all content and takes full responsibility for it. Separately, Claude Sonnet~5 is the evaluated advisor and thus part of the method (Section~\ref{sec:advisor}).

